% =====================================================================================
% sections/ext_discussion.tex: material cut from the Discussion and Limitations on 2026-10-10
% (CUT_PLAN.md, Sec. 4). Extended version only; appendix-ready (\subsection, no \section).
% Where each block came from (all from the pre-cut sections/discussion.tex unless noted):
%   "Status of the Measurements"       paragraph "Measurements", unchanged in substance
%   "Real Weights, Model Quality and V1" paragraphs "Real weights and model quality" and "V1";
%                                      "halves the proof" reworded per CUT_PLAN.md decision 4
%   "A Client's Options in Detail"     the long form of the body's "Options for a client"
%                                      paragraph and Table tab:options; numbers from
%                                      evaluation.tex (Secs. sec:eval-reexec, sec:eval-compare),
%                                      attack.tex (Table tab:sok), defence.tex (Sec.
%                                      sec:setup-proof), intro.tex and ext_intro.tex (gpucc,
%                                      CommitLLM). Derived by arithmetic only: ceil(11,008/2,048)
%                                      = 6 queries for the down projection.
%   "Limitations in Detail"            paragraphs "Proof size and the verifier", "Weights and
%                                      trust" and "The attack and the theorem" at full length
% The old Conclusion only summarised the body and is not kept here.
% Labels defined: ext:disc-status, ext:disc-real, ext:disc-options, ext:disc-limits (the body
%   does not reference them). Labels used (defined in the body or the full appendices):
%   tab:eval-llm, tab:eval-v1, tab:eval-reexec, tab:eval-compare, tab:options, tab:sok,
%   sec:eval-reexec, sec:setup-proof, sec:v1, sec:gen-sampling, sec:quality, rem:scope.
% =====================================================================================

\providecommand{\para}[1]{\par\noindent\textbf{#1}\quad}
\providecommand{\todo}[1]{\textbf{[TODO: #1]}}
\providecommand{\sC}{\mathsf{C}}
\providecommand{\sK}{\mathsf{K}}
\providecommand{\Kpre}{\mathsf{K}_{\mathrm{pre}}}

\subsection{Status of the Measurements}
\label{ext:disc-status}

The base implementation was measured on an NVIDIA L40S and an AMD EPYC 9334, and each improvement
and V1 on an RTX~2080~Ti and a Xeon Silver 4114, against the previous code on the same queries.
The improvements change no byte of the proof, so the proof sizes stand, but the times of the
frozen code on the L40S are not yet measured, and several cells of
Tables~\ref{tab:eval-llm}, \ref{tab:eval-v1} and~\ref{tab:eval-compare} are empty for that reason
\todo{L40S: final run of the frozen code; fill every \texttt{L40S} cell}. The 7B and 13B models
were measured as one to four of their blocks on the 2080~Ti, and full-model figures for them come
from the L40S run or from the byte model. Peak memory is not yet reported \todo{F4: peak memory}.

\subsection{Real Weights, Model Quality and V1}
\label{ext:disc-real}

\para{Real weights and model quality.} The costs were measured with random int8 weights. They
depend only on the shapes, except for the compressed size of the claims and of V1's values. On the
smoothed OPT-125M at 512 tokens, real activations compress less well than random ones as int8
values, and V1 shrinks the proof by 1.68$\times$ in mode $\sC$ (82.2 to 48.9\,MB) and
1.77$\times$ in $\Kpre$, less than with random weights (1.94--2.06$\times$ for GPT-2 at 512
tokens), also because seven of its residual operations have multipliers above $2^{29}$ and stay
clear under (P3) \todo{real-weight cells at 2{,}048 tokens; end-to-end run of Qwen3-4B and
OPT-6.7B (F3)}. Smoothing brings OPT-125M and OPT-1.3B within a factor of 1.03 of fp32
(Section~\ref{sec:quality}). We have not yet run it on OPT-6.7B, whose integer model with a scalar
gain is 33\% above fp32 (36.0 against 27.1), or on Qwen3-4B, because both need more memory than
our 2080~Ti setup provides. Soundness holds for any committed integer model; these numbers concern
only how well that model serves.

\para{V1.} V1 is an option, not a replacement. At 2{,}048 tokens it makes the proof
2.0--2.3$\times$ smaller with random weights, but it makes the prover about ten times slower on the
2080~Ti and the CPU verifier 1.6--1.9$\times$ slower. Llama-2-7B in mode $\sC$ does not fit on the
11\,GB card with the GKR, so V1 has not been timed in that configuration, and the full 7B and 13B
models have only the byte model's prediction \todo{L40S: V1 on the full Llama-2-7B, incl.\ mode
$\sC$, Llama-2-13B, OPT-6.7B and Qwen3-4B}. Mode $\sK$ and the streaming GPU verifier do not
support V1 yet, and V1 assumes one requantisation multiplier per operation, which excludes the
per-row weight scales that would bring OPT-125M to within $\times1.00$--$1.01$ of fp32.

\subsection{A Client's Options in Detail}
\label{ext:disc-options}

Table~\ref{tab:options} compares the options of a client that wants one Llama-2-7B answer checked.
We add the details behind each row.

\para{Re-execution.} A client that downloads the 6.7\,GB of int8 weights once can re-run every
query. The integer forward pass is bit-identical to the claims and so a complete check; it took
6.8\,s at 64 tokens and 196\,s at 2{,}048 tokens on 8 threads of a Xeon Silver 4114, and 0.16 and
2.08\,s on an RTX~2080~Ti, before the int8 forward pass that halves the GPU times
(Table~\ref{tab:eval-reexec}). Floating-point re-execution is faster (fp32 on the CPU: 3.5 and
86\,s; fp16 on the GPU: 0.037 and 0.57\,s) but cannot confirm an exact answer. A CPU client that
holds the weights can instead run our $\Kpre$ verifier, which takes 1.5 and 85\,s, 2.3--4.5$\times$
faster than re-execution at 64 tokens and 2.3$\times$ faster than integer re-execution at 2{,}048.
On a GPU, re-execution wins: the base GPU verifier needed 4.95\,s in $\Kpre$ at 2{,}048 tokens on
the faster L40S (Section~\ref{sec:eval-reexec}).

\para{GPU confidential computing.} An H100-class GPU runs the whole model inside an attested
enclave. A measurement study finds little overhead inside the GPU and attributes most of the
penalty, below 7\% for most LLM queries and close to zero for larger models and longer sequences,
to encrypted CPU--GPU transfers over PCIe~\cite{gpucc}. The client checks an attestation rather
than a proof, and the weights stay hidden from it. The guarantee rests on the vendor's hardware,
its attestation service and the enclave's resistance to side channels.

\para{CommitLLM-style audit.} CommitLLM~\cite{commitllm2026} adds 12--14\% to serving. Its routine
audit opens one token of 1--5\% of the responses and 10 of 80 layers at that token, so an edit
passes with probability 87.5\% when the token is opened and about 99.4\% per one-token response at
a 5\% audit rate (Table~\ref{tab:sok}). Its client holds secret Freivalds vectors precomputed from
the weights, and attention is audited, not verified \todo{CommitLLM's bytes and client time per
audited response}.

\para{zkSNARKs.} zkLLM~\cite{zkllm} proves a 2{,}048-token Llama-2-7B prompt in 10.3\,min on an
A100, with a 183\,kB proof that its verifier checks in 2.36\,s, and hides the weights. zkAgent,
which proves a generated sequence in one pass, reports 0.30\,s per token for GPT-2~\cite{zkagent}.
On a 2{,}048-token Llama-2-13B prompt our L40S prover is 46$\times$ faster than zkLLM, but zkLLM
also proves the operations without weights and hides the weights (Table~\ref{tab:eval-compare}).

\para{Ours.} In mode $\sC$ the client holds a 32-byte root per matrix and no weights. A 64-token
proof is 325\,MB, about 21$\times$ smaller than the int8 weights, and is checked in 1.0\,s on the
EPYC CPU or 431\,ms on the L40S. Of it, 124.7\,MB are opened columns, which do not shrink with
shorter prompts. At 2{,}048 tokens the proof is 6.6\,GB (6.4\,GB in $\Kpre$), about the size of the
weights, and takes 5.3\,s to transfer at 10\,Gb/s \todo{F7: transfer times}. A client that sends
more than about 21 short prompts, or one long one, therefore receives more bytes than a download of
the weights, and the protocol pays off only for a client that cannot store or run the model, or
will not. Our base prover took 1.8\,s for a 64-token prompt on the L40S, against 0.037\,s for an
fp16 forward pass on the 2080~Ti, on different GPUs \todo{L40S: prover against unverified fp16
and int8 inference, same GPU}. Sending one opening and one fold for a batch of queries whose
claims are all fixed before the challenges would amortise the opened columns over the batch; we
have not implemented this.

\subsection{Limitations in Detail}
\label{ext:disc-limits}

\para{Proof size and the verifier.} Even with V1, the proof grows with the prompt, because it
carries one value for every output of every weight operation. For Llama-2-7B it reaches the size of
the int8 weights at about 2{,}050 tokens without V1 and, by the byte model, at about 4{,}500 tokens
with it. Only committing the activations and proving the operations without weights, as zkLLM
does, would remove this growth; we leave it to future work. The verifier recomputes attention, so
its cost grows with the square of the prompt length; at 2{,}048 tokens the operations without
weights take 76\% of the CPU verifier's time \todo{B5: native attention kernel}. A CPU client that
holds the weights saves time over re-execution at 64 tokens and ties with fp32 re-execution at
2{,}048 tokens, and on a GPU re-execution is faster than our verifier. The case for the protocol is
therefore a client that does not hold the weights or the hardware to run them, not a faster way to
compute an answer.

\para{Weights and trust.} The protocol is not zero-knowledge: the claims of about
$\lceil k/T\rceil$ queries reveal a weight matrix with rows of length $k$ (about six 2{,}048-token
queries for Llama-2-7B's down projection, $k=11{,}008$), so mode $\sC$ targets open-weight models.
For private weights it shares the limitation that Hollow-LLM points out for
zkSNARKs~\cite{hollowllm}: the commitment binds the provider to one fixed model, but a client that
sees neither the weights nor a range proof on them does not learn which model that is, and weights
with a degenerate structure keep the declared shapes while being cheap to evaluate. With an
untrusted commitment (Section~\ref{sec:setup-proof}), the bound (P4) of V1 and the claim range of
v0 must be published with the model, and the time of the setup proof is not yet measured
\todo{L40S: setup prover and verifier time}. The exact sampler (Section~\ref{sec:gen-sampling}) is
a fixed approximation of softmax sampling whose distance to it we have not measured, the proof
does not show that the client's seed was random, and nucleus sampling is not implemented.

\para{The attack and the theorem.} The lower bound needs condition (H), which we establish for
random ensembles with independent rows and on a constructed worst-case network, not for a trained
network, whose rows are correlated. On trained networks we report the detection of the best attacks
we found, which bounds each sampler's worst-case detection from above. We do not claim that the
worst-case instance embeds in a transformer, since the LayerNorm before the MLP prevents feeding it
the required patterns (Remark~\ref{rem:scope}). The natural ranges in our range-respecting attacks
were computed on test queries that include the attacked ones, and the adaptive game against
value-aware samplers with held-out ranges is not yet run \todo{G3: adaptive game, held-out
ranges}. Our OPT-6.7B attack writes values outside the neurons' natural ranges; an in-range attack
on an LLM is pending \todo{G3.5: in-range LLM attack}. We deliberately do not develop a method for
steering the responses of instruction-tuned assistants.
